\documentclass[11pt]{article}

\usepackage{amsmath,amssymb}
\usepackage{xurl}
\usepackage[hidelinks]{hyperref}
\setlength{\emergencystretch}{2em}

% Shared commands for notes in docs/paper-gaps/.

\DeclareUnicodeCharacter{2080}{\ifmmode{}_{0}\else\textsubscript{0}\fi}
\DeclareUnicodeCharacter{2081}{\ifmmode{}_{1}\else\textsubscript{1}\fi}
\DeclareUnicodeCharacter{2082}{\ifmmode{}_{2}\else\textsubscript{2}\fi}
\DeclareUnicodeCharacter{2099}{\ifmmode{}_{n}\else\textsubscript{n}\fi}

\newcommand{\Lean}{\textsc{Lean}}
\ExplSyntaxOn
\NewDocumentCommand{\leanid}{m}{
  \ifmmode
    \textsf{\detokenize{#1}}
  \else
    \group_begin:
    \urlstyle{sf}
    \tl_set:Nn \l_tmpa_tl {#1}
    \tl_replace_all:Nnn \l_tmpa_tl {\_} {_}
    \exp_args:NV \path \l_tmpa_tl
    \group_end:
  \fi
}
\ExplSyntaxOff
\providecommand{\tr}{\operatorname{tr}}
\newcommand{\tnleanIssueBase}{https://github.com/LionSR/TNLean/issues/}
\newcommand{\tnleanPrBase}{https://github.com/LionSR/TNLean/pull/}
\newcommand{\tnleanDocBase}{https://sirui-lu.com/TNLean/docs/}
\newcommand{\ghissue}[1]{\href{\tnleanIssueBase#1}{GitHub issue \##1}}
\newcommand{\ghpr}[1]{\href{\tnleanPrBase#1}{PR \##1}}
\newcommand{\doclink}[1]{\href{\tnleanDocBase#1.html}{\path{#1}}}

% Dirac notation and the identity operator, as in blueprint/src/macros/common.tex.
% \providecommand keeps note-local definitions authoritative.
\usepackage{dsfont}
\providecommand{\ket}[1]{|#1\rangle}
\providecommand{\bra}[1]{\langle #1|}
\providecommand{\braket}[2]{\langle #1 | #2 \rangle}
\providecommand{\ketbra}[2]{|#1\rangle\!\langle #2|}
\providecommand{\Id}{\mathds{1}}
\providecommand{\one}{\mathds{1}}
\providecommand{\C}{\mathbb{C}}
\providecommand{\MN}[1]{M_{#1}(\C)}
\providecommand{\MD}{M_D(\C)}

% External referencing for paper sources.  The build helper writes sanitized
% auxiliary files under build/paper-aux/ and excludes duplicated source labels.
\usepackage{xr}

% Import a generated paper auxiliary file with a caller-supplied label prefix.
% The candidate paths support builds from docs/paper-gaps/, the repository root,
% and build/paper-aux/ itself.
\newcommand{\importpaperaux}[2]{%
  \IfFileExists{../../build/paper-aux/#2.aux}{%
    \externaldocument[#1]{../../build/paper-aux/#2}%
  }{%
    \IfFileExists{build/paper-aux/#2.aux}{%
      \externaldocument[#1]{build/paper-aux/#2}%
    }{%
      \IfFileExists{#2.aux}{%
        \externaldocument[#1]{#2}%
      }{}%
    }%
  }%
}

% General external reference macros
\newcommand{\paperref}[2]{\ref{#1-#2}}
\newcommand{\papereqref}[2]{(\ref{#1-#2})}

% CPSV16 source (1606.00608 / MPDO-22-12-17-2.tex) external document and shortcuts
\importpaperaux{cpsv16-}{1606.00608/MPDO-22-12-17-2-tnlean}
\newcommand{\cpsvref}[1]{\paperref{cpsv16}{#1}}
\newcommand{\cpsveqref}[1]{\papereqref{cpsv16}{#1}}

% Verdict marker: \gapnote{<kind>}{<status>}. Kind is the policy
% classification (clarification, local-correction, scope-restriction,
% unfaithful, false-source, open-gap); status is open, wip (elimination
% underway), resolved, or historical. Machine-read by the site index;
% prints a verdict line.
\newcommand{\gapnote}[2]{\par\noindent\textbf{Verdict:} #1 (#2).\par}


\title{The Root-Channel Kraus-Rank Step in Theorem 4.1}
\author{}
\date{2026-04-30}

\begin{document}
\maketitle

\section{The assertion}

The reverse implication in Theorem~4.1 of de las Cuevas, Cirac, Schuch,
and P\'erez-Garc\'ia is the point at issue
\cite[Theorem~4.1]{DeLasCuevas2017Irreducible}.\footnote{For traceability, this
note corresponds to \href{https://github.com/LionSR/TNLean/issues/1046}{issue \#1046}.
Related formal work is recorded in
\href{https://github.com/LionSR/TNLean/issues/670}{issue \#670},
\href{https://github.com/LionSR/TNLean/issues/821}{issue \#821},
\href{https://github.com/LionSR/TNLean/issues/622}{issue \#622}, and
\href{https://github.com/LionSR/TNLean/issues/664}{issue \#664}.} The source passages
are \path{Papers/1708.00029/main.tex:717--731} and
\path{Papers/1708.00029/main.tex:812--818}.

Let $B=(B^1,\ldots,B^d)$ be an MPS tensor, let $p\geq 1$, and let
$\mathcal E_B$ be the transfer channel
\[
  \mathcal E_B(X)=\sum_{i=1}^d B^i X (B^i)^\dagger.
\]
The paper says that a trace-preserving completely positive map $\mathcal E$ is
$p$-divisible if there is another trace-preserving completely positive map
$\mathcal E'$ with
\[
  \mathcal E=(\mathcal E')^p.
\]
This is a channel-level condition. It places no restriction on the number of
Kraus operators needed to represent the root channel $\mathcal E'$.

In the same passage, $B$ is called $p$-refinable if there are a tensor
$A=(A^1,\ldots,A^d)$ with the same physical alphabet and an isometry
\[
  W:\mathbb C^d \longrightarrow (\mathbb C^d)^{\otimes p}
\]
such that
\[
  |V_{pN}(A)\rangle = W^{\otimes N}|V_N(B)\rangle
  \qquad\text{for every } N.
\]
Thus the tensor root in a refinement is indexed by the original $d$ physical
letters. The theorem states the unconditional equivalence
\[
  B \text{ is } p\text{-refinable}
  \quad\Longleftrightarrow\quad
  \mathcal E_B \text{ is } p\text{-divisible}
\]
for $B$ in irreducible form~II.

\section{The point at issue}

The reverse proof begins with $p$-divisibility of $\mathcal E_B$. This gives a
trace-preserving completely positive map $\mathcal E'$ such that
$\mathcal E_B=(\mathcal E')^p$. The proof then writes that there exists a tensor
$A$ such that
\[
  \mathcal E_B=\mathcal E_A^p,
\]
and proceeds to compare two Kraus representations of the same channel by the
usual isometric freedom of Kraus representations.

The displayed passage is the missing step. To realize the channel root
$\mathcal E'$ as $\mathcal E_A$ for a tensor
$A=(A^1,\ldots,A^d)$, the root channel must admit a Kraus representation with at
most $d$ Kraus operators. Equivalently, its minimal Kraus rank, or Choi rank,
must be at most $d$. If the root channel has minimal Kraus rank $r>d$, then it
can be represented by an $r$-letter tensor but not by a $d$-letter tensor. The
resulting isometry would land in $(\mathbb C^r)^{\otimes p}$ rather than in the
space $(\mathbb C^d)^{\otimes p}$ required by the definition of $p$-refinement.

This is a mathematical point, not a matter of notation. A completely positive map
can require more than a prescribed number of Kraus operators. The channel-level
assertion $\mathcal E_B=(\mathcal E')^p$ does not by itself rule out the possibility that
every Kraus representation of $\mathcal E'$ has more than $d$ terms. Therefore
one needs either an argument that the special hypotheses of Theorem~4.1 force a
$d$-term Kraus representation for some root channel, or a different proof of the
reverse implication.

\section{The formal statement}

The formal definition of $p$-divisibility follows the channel-level definition
from the paper: it asserts the existence of a trace-preserving completely
positive root channel and does not include any bound on the root channel's Kraus
rank. The conditional reverse theorem therefore separates the missing step from
the rest of the proof.
\footnote{The definition is
\texttt{MPSTensor.IsPDivisibleChannel} in
\path{TNLean/MPS/Periodic/Symmetry/Theorem41Defs.lean:24--32}. The conditional
reverse theorem is \texttt{MPSTensor.thm\_4\_1\_p\_refinement\_reverse} in
\path{TNLean/MPS/Periodic/Symmetry/Theorem41Reverse.lean:120--142}.}

The extra formal hypothesis is an inverse canonicalization statement. In
ordinary mathematical terms, it says that whenever $B$ is in irreducible form~II
and $\mathcal E_B$ is $p$-divisible, there is a $d$-letter tensor $A$ such that
\[
  \mathcal E_B=\mathcal E_{A^{[p]}}.
\]
If such a tensor-level root is available, the rest of the argument follows
the source proof: the tensor $A^{[p]}$ obtained by blocking and the original
tensor $B$ give two finite Kraus representations of the same completely positive map, and
isometric freedom of Kraus representations gives the isometry
$W:\mathbb C^d\to(\mathbb C^d)^{\otimes p}$ needed for refinement.
\footnote{The inverse canonicalization hypothesis is
\texttt{MPSTensor.PRefinementInverseCanonicalization} in
\path{TNLean/MPS/Periodic/Symmetry/Theorem41Reverse.lean:26--42}. The blueprint
records the corresponding conditional statement in
\path{blueprint/src/chapter/ch11b_periodic_ft.tex:1076--1103}.}

The formal theorem also isolates a sufficient channel-theoretic hypothesis for
the missing step: if every trace-preserving completely positive root
$\mathcal E'$ with $\mathcal E_B=(\mathcal E')^p$ has Kraus rank at most $d$,
then one can choose a $d$-term Kraus representation by padding with zero Kraus
operators and obtain the required tensor $A$. This is a genuine mathematical
hypothesis about the root channel, not a change of notation. The minimal-Kraus
fact that the least number of Kraus operators equals the Choi rank is the right
background theorem, but it does not itself prove the required inequality
$\operatorname{rank}_{\mathrm{Kraus}}(\mathcal E')\le d$ for a root of
$\mathcal E_B$.
\footnote{The bounded-root formulation is
\texttt{MPSTensor.PRefinementInverseRootKrausRankBound}; the implication from
this bound to inverse canonicalization is
\texttt{MPSTensor.pRefinementInverseCanonicalization\_of\_rootKrausRankBound},
recorded in \path{TNLean/MPS/Periodic/Symmetry/Theorem41Reverse.lean:98--118}.}

\section{Conclusion}

The formal theorem adds a hypothesis relative to the source theorem. The source
statement claims that channel-level $p$-divisibility of $\mathcal E_B$ is already
enough to construct a $p$-refinement of $B$, but the reverse proof passes through
a tensor root whose existence requires a Kraus-rank bound on the channel root.
That bound is not part of the paper's definition of $p$-divisibility and is not
proved in the cited converse paragraph.

The point left open is whether the channel root can always be chosen with Kraus
rank at most $d$ under the irreducible-form hypotheses of the theorem. The formal
theorem therefore proves the reverse implication under an explicit inverse
canonicalization hypothesis, or under the stronger root-Kraus-rank assumption
which implies it. Until a paper-level argument proves such a bound, the reverse
implication of Theorem~4.1 should be treated in the formal statement as
conditional on that additional mathematical input.

\bibliographystyle{alpha}
\bibliography{references}

\end{document}